\documentclass[11pt]{article}
\usepackage[T1]{fontenc}
\usepackage[a4paper,left=2.1cm,right=2.1cm,top=1.8cm,bottom=1.8cm]{geometry}
\usepackage{amsmath,amssymb,booktabs,tabularx,array,hyperref}
\usepackage{algorithm,algpseudocode}
\hypersetup{hidelinks}
\setlength{\emergencystretch}{3em}
\newcommand{\Nfit}{N_{\mathrm{fit}}}
\newcommand{\Sstore}{S_{\mathrm{store}}}
\newcommand{\Sfit}{S_{\mathrm{fit}}}
\newcommand{\Sq}{S_q}
\newcommand{\Kmax}{K_{\max}}
\newcommand{\eps}{\varepsilon}
\newcommand{\code}[1]{\texttt{\detokenize{#1}}}
\DeclareMathOperator{\MSE}{MSE}
\DeclareMathOperator{\MAE}{MAE}
\title{Fully Online Retrieval Adaptation\\Experiment Guideline}
\author{}
\date{29 August 2026}

\begin{document}
\maketitle

\section{Problem}

Let $Z=(z_{t,u})\in\mathbb R^{T\times U}$ be one time-series panel, where
$t\in\{0,\ldots,T-1\}$ is a date index and
$u\in\{1,\ldots,U\}$ is a series (called a user in the code). For a query
date $q$, a lookback length $L$, and a forecast horizon $H$, define
\[
 x_{q,u}=(z_{q-L+1,u},\ldots,z_{q,u})\in\mathbb R^L,
 \qquad
 y_{q,u}=(z_{q+1,u},\ldots,z_{q+H,u})\in\mathbb R^H.
\]
The task is to predict every $y_{q,u}$. A precomputed date plan separates the
T0 datastore interval, T1+T2 fitting interval, and T3 evaluation grid. At date
$q$, a row dated $r$ may be retrieved or fitted only if
\[
 r+H\le q,                                                     \tag{1}
\]
so its complete target $y_{r,u}$ is already observed. The experiment asks
whether these causal rows improve a frozen forecasting model.

The CSV loader discovers \code{config.json} beside the selected dataset.
Portable top-level values are overridden by the \code{online_adaptation}
object and then by explicit run settings. For \code{drop_users}, null or
omission inherits the preceding level, an empty list retains every CSV user,
and a nonempty value replaces the preceding list. Exclusions are applied once
while reading the CSV and the effective list is recorded in the run manifest.
After aggregation, \code{missing_values=zero} replaces every NaN with zero;
\code{missing_values=error} rejects the dataset. Infinite values are never
repaired by either policy.

\section{Mathematical definition}

\subsection{Datastore, fitting dates, and queries}

Let $P$ be the dataset period: $P=24$ for hourly data, $P=7$ for daily data,
and $P=96$ for 15-minute data. Let $\Sstore$, $\Sfit$, and $\Sq$ be the
datastore, fitting, and evaluation-query strides. With periodic alignment,
the latest causal datastore date for query $q$ is
\[
 b_q=q-P\left\lceil\frac HP\right\rceil .                    \tag{2}
\]
Without alignment, $b_q=q-H$. The available datastore-date grid is
\[
 \mathcal D_q^\infty
 =\{b_q-j\Sstore:j\in\mathbb N_0,\ b_q-j\Sstore\ge L-1\},    \tag{3}
\]
where $\mathbb N_0=\{0,1,2,\ldots\}$. Aligned runs require $\Sstore$ to be a
multiple of $P$.

Let $M$ be \code{N\_datastore\_dates}: a positive integer number of retained
window origins, or the floor of a configured ratio times the number of complete
origins on the datastore-stride grid. Scope does not change $M$: one same-user candidate, $U$ all-user
candidates, or $U-1$ other-user candidates are available per retained date.
The precomputed T0 interval is
\[
 [a_0,b_0]=[L-1,\ L-1+M\Sstore-1].                            \tag{4}
\]
A rolling $\mathcal D_q$ contains the latest $M$ dates of
$\mathcal D_q^\infty$. A fixed datastore uses the same immutable T0 interval
for every query and subsets it by the query's periodic phase when alignment is
enabled. The first retrieval window occurs only after the full T0 capacity and
all its horizons are observable; partial early datastores are not evaluated.

Fitting dates are selected independently of $\mathcal D_q$. Let
\[
 c_q=
 \begin{cases}
 q-P\lceil H/P\rceil,
   & \text{alignment on and }P\mid\Sfit,\\
 q-H, & \text{otherwise},
 \end{cases}
\]
where $P\mid\Sfit$ means that $\Sfit$ is a multiple of $P$. Let $r_0$ be the
first retrieval window after T0. The rolling fitting set is the latest
$\Nfit$ dates in
\[
 \mathcal F_q^\infty
 =\{c_q-j\Sfit:j\in\mathbb N_0,\ c_q-j\Sfit\ge r_0\}.        \tag{5}
\]
Here $\Nfit$ always counts dates per user. Same-user fitting uses $\Nfit$
rows for each user; all-user fitting uses $\Nfit U$ rows for one shared fit.
A fixed training set is the single precomputed T1+T2 grid
$\{r_0+j\Sfit:j=0,\ldots,\Nfit-1\}$ and produces one ridge reused for all
evaluation queries.

The natural first evaluated date $q_0$ follows the complete T0 interval and
all $\Nfit$ T1+T2 origins after their horizons become observable. An explicit
absolute or fractional evaluation start may move $q_0$ later, never earlier;
an optional end boundary may shorten T3. The evaluated dates are fixed before
extraction as
\[
 \mathcal Q=\{q_0+j\Sq:j\in\mathbb N_0,\ q_0+j\Sq\le T-H-1\}.\tag{6}
\]

\subsection{Retrieval}

Let $\phi(x)\in\mathbb R^J$ be the retrieval representation of a lookback,
where $J$ is its number of features, and let $d$ be the configured distance.
The default is $\phi(x)=x$ and
\[
 d(x,x')=\|x-x'\|_2.                                         \tag{7}
\]
For each $(q,u)$, exact search returns the $\Kmax$ smallest-distance windows
$(r_{q,u,k},v_{q,u,k})$, $k\in\{1,\ldots,\Kmax\}$, where
$r_{q,u,k}\in\mathcal D_q$ is a date and $v_{q,u,k}$ is a user. The allowed
users are all $v$, only $v=u$, or only $v\ne u$, according to retrieval scope.

Let $\mu(x)$ and $\sigma(x)$ be the mean and population standard deviation of
the entries of $x$, and let $\eps=10^{-8}$. A value $a$ attached to a neighbor
lookback $x'$ is put on the query scale by
\[
 R_x(a;x')=
 \frac{a-\mu(x')}{\max\{\sigma(x'),\eps\}}
 \max\{\sigma(x),\eps\}+\mu(x).                             \tag{8}
\]

\subsection{Signals used by the ridge}

Let $f:\mathbb R^L\to\mathbb R^H$ be the frozen forecaster. For every
$(q,u)$ and neighbor rank $k$, define
\begin{align*}
 V_{q,u} &= f(x_{q,u}), &&\text{vanilla forecast},\\
 Y_{q,u,k} &= R_{x_{q,u}}(y_{r_{q,u,k},v_{q,u,k}};
                         x_{r_{q,u,k},v_{q,u,k}}),
 &&\text{retrieved future},\\
 N_{q,u,k} &= R_{x_{q,u}}(f(x_{r_{q,u,k},v_{q,u,k}});
                         x_{r_{q,u,k},v_{q,u,k}}),
 &&\text{forecast of the retrieved lookback}.
\end{align*}
For a chosen $K\le\Kmax$, the distance-weighted retrieved future is
\[
 \bar Y_{q,u,K}=\sum_{k=1}^K w_{q,u,k}Y_{q,u,k},\qquad
 w_{q,u,k}=\frac{e^{-\delta_{q,u,k}}}
                  {\sum_{j=1}^K e^{-\delta_{q,u,j}}},
\]
\[
 \delta_{q,u,k}=\frac{d_{q,u,k}-\min_{j\le K}d_{q,u,j}}
 {\max\{\operatorname{sd}(d_{q,u,1:K}),\eps\}},             \tag{9}
\]
where $d_{q,u,k}$ is the stored retrieval distance and $\operatorname{sd}$ is
population standard deviation.

The code also defines a context forecast $C_{q,u,K}$ by calling the canonical
backbone with the first $K$ rescaled neighbor pairs translated into named
covariates. With \path{retrieval_covariate_mode=past_and_future}, the past
covariates are the neighbor lookbacks and the future covariates are their known
horizons; \code{past} omits the latter and \code{none} disables the wrapper.
Chronos-2 therefore computes
\[
 C_{q,u,K}=f\!\left(x_{q,u};\widetilde X_{q,u,1:K},
                    \widetilde Y_{q,u,1:K}\right),             \tag{10}
\]
through the official pipeline. The model alias remains \code{chronos2} in all
three modes. Passing context to a backbone that declares no covariate support
raises an error; it is never replaced by $V_{q,u}$. The four-backbone profile
uses the covariate-free \code{y_ridge_shared} design, while the primary
Chronos-2 ridge and gate profiles use the context forecast.

For one date, user, and horizon position, the feature vector is selected from
the following signals. Each $Y$ or $N$ entry denotes all ranks $1,\ldots,K$.
\begin{center}
\small
\begin{tabular}{@{}ll@{}}
\toprule
Design name & Feature signals \\
\midrule
\texttt{cov} & $V,C$ \\
\texttt{avgy} & $V,\bar Y$ \\
\texttt{y} & $V,Y_1,\ldots,Y_K$ \\
\texttt{cov\_y} & $V,C,Y_1,\ldots,Y_K$ \\
\texttt{cov\_avgy} & $V,C,\bar Y$ \\
\texttt{residual} & $V,Y_1,\ldots,Y_K,N_1,\ldots,N_K$ \\
\texttt{full} & $V,C,Y_1,\ldots,Y_K,N_1,\ldots,N_K$ \\
\bottomrule
\end{tabular}
\end{center}

The main method is \texttt{full\_ridge\_shared}. Let $D\in\mathbb R^{nH\times p}$
contain its $p=2+2K$ features for $n$ fitting date--user rows, and let
$e\in\mathbb R^{nH}$ contain $y-V$. If the fitting loss is nMSE, each date--user
row of $D$ and $e$ is first divided by
$\max\{\sigma(x_{r,u}),\eps\}$; for MSE it is divided by $1$. Let
\[
 S=\operatorname{diag}(s_1,\ldots,s_p),\qquad
 s_j=\max\left\{\sqrt{\frac1{nH}\sum_{i=1}^{nH}D_{ij}^2},10^{-12}\right\}.
\]
For ridge value $\alpha\ge0$, the fitted coefficients are
\[
 \widehat\beta
 =S^{-1}\left(\frac{S^{-1}D^\top D S^{-1}}{nH}+\alpha I_p\right)^{-1}
 \frac{S^{-1}D^\top e}{nH},                                  \tag{11}
\]
where $I_p$ is the $p\times p$ identity matrix. The prediction is
\[
 \widehat y_{q,u}=V_{q,u}+D_{q,u}\widehat\beta.              \tag{12}
\]
Same-user fitting gives one $\widehat\beta$ per $(q,u)$; all-user fitting gives
one $\widehat\beta$ per $q$. Shared ridge uses one coefficient vector for all
$H$ forecast positions. Horizon ridge fits one vector $\widehat\beta_h$ for
each $h\in\{1,\ldots,H\}$.

Delta ridge removes $V$ from $D$, replaces every remaining signal $A$ by
$A-V$, and still predicts (12). Convex fitting uses the unmodified signals,
targets $y$, and solves
\[
 \min_{\beta\in\mathbb R^p}\frac1{nH}\|y-D\beta\|_2^2
 \quad\text{subject to}\quad \beta_j\ge0,
 \quad\sum_{j=1}^p\beta_j=1;                                \tag{13}
\]
its prediction is $D\widehat\beta$.

\subsection{Selection at each query}

Let $\rho\in(0,1)$ be the validation ratio and
$m=\lceil\rho\Nfit\rceil$. When $\alpha$ or $K$ is selected, the oldest
$\Nfit-m$ dates of $\mathcal F_q$ fit each candidate and the newest $m$ dates
score it. The default candidate sets are
\[
 \mathcal A=\{10^{-1},10^{-2},10^{-3}\},\qquad
 \mathcal K=\{1,5,10,15\},\qquad \Kmax=20,\qquad \rho=0.2. \tag{14}
\]
The selected pair minimizes validation MSE or nMSE. Exact ties choose smaller
$K$, then larger $\alpha$. The selected pair is refitted on all of
$\mathcal F_q$. Supplying \texttt{used\_k} fixes $K$ but does not fix
$\alpha$; setting \texttt{tune\_alpha=False} fixes $\alpha$ but does not fix
$K$.

\subsection{Gate ablations and TS-RAG comparison}

For a gate candidate $A_{r,u}\in\{C_{r,u,K},\bar Y_{r,u,K}\}$, define its
horizon-averaged advantage
\[
 a_{r,u}=\frac1H\sum_{h=1}^H
 \frac{(y_{r,u,h}-V_{r,u,h})^2-(y_{r,u,h}-A_{r,u,h})^2}
 {c_{r,u}^2},                                                 \tag{15}
\]
where $c_{r,u}=1$ for MSE fitting and
$c_{r,u}=\max\{\sigma(x_{r,u}),\eps\}$ for nMSE fitting. For the method named
\texttt{bayes}, the score $s_{q,u}$ is the mean of $a_{r,u}$ over the fitting
rows. For \texttt{catboost}, define $\Delta=A-V$,
$E=Y_{1:K}-V$, $G=Y_{1:K}-N_{1:K}$, the rescaled neighbor lookbacks
$X_k=R_{x_{r,u}}(x_{r_k,v_k};x_{r_k,v_k})$, neighbor ages
$g_k=r-r_k$, and same-user indicators $o_k=\mathbf 1[v_k=u]$. Means and
population standard deviations below use every entry of the indicated array.
The regressor input is
\begin{align*}
 \psi_{r,u}=(&\operatorname{mean}\Delta,\operatorname{sd}\Delta,
 \operatorname{mean}E,\operatorname{sd}E,
 \operatorname{mean}G,\operatorname{sd}G,
 \mu(x_{r,u}),\sigma(x_{r,u}),\\
 &\operatorname{mean}X,\operatorname{sd}X,
 \operatorname{mean}o,\operatorname{mean}g,\operatorname{sd}g,
 \operatorname{mean}d,\operatorname{sd}d,\min d)\in\mathbb R^{16}. \tag{16}
\end{align*}
CatBoost is fitted to map $\psi_{r,u}$ to $a_{r,u}$, and its output is
$s_{q,u}$. Let
$\tau_{q,u}$ be the population standard deviation of the fitting advantages.
Both methods predict
\[
 \widehat y_{q,u}=V_{q,u}+
 \operatorname{sigmoid}(s_{q,u}/\tau_{q,u})(A_{q,u}-V_{q,u}),\tag{17}
\]
where $\operatorname{sigmoid}(x)=(1+e^{-x})^{-1}$. When $\tau_{q,u}=0$, the
weight is $1$, $0$, or $1/2$ for $s_{q,u}>0$, $s_{q,u}<0$, or $s_{q,u}=0$.
The primary workflow uses $A=C$ and $K=20$ for its Bayes baseline.

The TS-RAG comparison is source-adapted from
\href{https://github.com/UConn-DSIS/TS-RAG/tree/73ac807789d2e61b8a3dfc8514e3fc947fe185cc}
{UConn-DSIS/TS-RAG commit 73ac807}. It is fixed to $L=512$, $H=64$, $K=5$,
Chronos-Bolt, and the released \code{only\_self\_train} scope: each variable
retrieves only from its own immutable training split. Retrieval origins are
unstrided and are not periodically aligned. Its retriever uses the final Chronos-T5
token $\phi(x)$, float32 \texttt{faiss.IndexFlatL2}, a $K+1$ search, and removal
of the final returned column. Its ARM receives the original, unscaled retrieved
lookback--future sequences.

\section{Measures}

For prediction $\widehat y_{q,u}$ and query scale
$s_{q,u}=\max\{\sigma(x_{q,u}),\eps\}$, define
\[
 \MSE_{q,u}=\frac1H\sum_{h=1}^H
 (y_{q,u,h}-\widehat y_{q,u,h})^2,
 \qquad
 \mathrm{nMSE}_{q,u}=\frac{\MSE_{q,u}}{s_{q,u}^2},            \tag{18}
\]
\[
 \MAE_{q,u}=\frac1H\sum_{h=1}^H
 |y_{q,u,h}-\widehat y_{q,u,h}|,
 \qquad
 \mathrm{nMAE}_{q,u}=\frac{\MAE_{q,u}}{s_{q,u}}.             \tag{19}
\]
Every aggregate is an unweighted mean over the retained date--user rows.
Compared methods must have identical query dates; a mismatch is an error. If $m$ denotes an
adapted method and $v$ its matching vanilla forecast, report
\[
 I^{\MSE}_m=100\frac{\MSE_v-\MSE_m}{\MSE_v},\qquad
 I^{\mathrm{nMSE}}_m=100\frac{\mathrm{nMSE}_v-\mathrm{nMSE}_m}
 {\mathrm{nMSE}_v},                                          \tag{20}
\]
and the strict win rate
\[
 W_m=100\frac1{|\mathcal Q|U}
 \sum_{q\in\mathcal Q}\sum_{u=1}^U
 \mathbf 1[\MSE_{m,q,u}<\MSE_{v,q,u}],                       \tag{21}
\]
where $\mathbf 1[\cdot]$ is $1$ when its condition is true and $0$ otherwise.
The report also saves population standard deviation across per-user means,
the mean over the worst $\lceil0.1U\rceil$ users selected separately for MSE
and nMSE, and the percentage of users improved.

Reported time is extraction time plus adaptation/evaluation time. It is also
divided by $|\mathcal Q|U$. A separate cold-batch value measures the first
all-user query with empty in-memory caches.

\section{Runs}

The default ridge run uses raw Euclidean all-user retrieval, same-user fitting,
non-fixed datastore and training set, $M=100$, $\Nfit=10$,
$\Sstore=\Sfit=P$, seed $1$, and (14). Test runs use $\Sq=257$; small and full
runs use $\Sq=127$. The frozen backbone is Chronos-2. The main family compares
full ridge, the vanilla-plus-retrieved-outcome convex mixture, the Bayes
covariate gate, direct covariate prediction, and vanilla. The separate TS-RAG
family uses only $512{:}64$ and the configuration defined above.

For the deadline experiment, four Selena-only fronts narrow this grid. The fixed
protocol uses Electricity and Solar only at the hourly short setting
$(L,H)=(168,24)$. It fixes T0/T1+T2/T3 to 30/50/20 percent of each source
timeline, uses periodically aligned datastore origins, fit stride 24, and T3
query stride 127. Every arm uses exactly $\Nfit=10$ fitting dates. The full
$2\times2\times2$ comparison crosses fixed versus rolling
datastore, fixed versus online fitting set, and same-user versus all-user
fitting. The first two fronts isolate the fully online same-user arm and the
fully fixed all-user arm; a dependent front runs the other six arms. Every arm
shares the exact T3 query-date grid and a global maximum of 10,000 datastore
windows. Ridge selection uses the chronologically newest 20 percent of its
integer fitting-date set for validation.

The priority TS-RAG comparison fixes $(L,H)=(512,64)$ for Electricity, Solar,
ETTh1, ETTh2, ETTm1, ETTm2, \path{ne_china_wind_h},
\path{coastal_t_s_h_part11}, and \path{sg_weather_d}; Traffic and all remaining
TIME panels stay in the main full comparison. Both methods evaluate the same
stride-127 T3 grid beginning at the 80-percent boundary. Full online ridge
first anchors T3 and 10 fitting dates at stride 24, then fills the preceding
datastore up to 10,000 windows. Its datastore is periodically aligned with
stride 24 for hourly data, 96 for 15-minute data, and 7 for daily data, and it
selects $K$ from $\{1,5,10,15\}$ together with alpha. TS-RAG instead uses its
fixed, unstrided, same-variable T0 datastore, the same global 10,000-window
maximum, and $K=5$.

\begin{center}
\small
\begin{tabularx}{\linewidth}{@{}p{.14\linewidth}X@{}}
\toprule
Mode & Datasets and $(L,H)$ settings \\
\midrule
test & Electricity, long range only \\
small & Electricity, Solar, and Traffic; short, mid, and long ranges \\
full & small plus raw ETTh1, ETTh2, ETTm1, ETTm2 and every prepared TIME
dataset; short, mid, and long ranges \\
\bottomrule
\end{tabularx}
\end{center}

\begin{center}
\small
\begin{tabular}{@{}lccc@{}}
\toprule
Frequency & short & mid & long \\
\midrule
hourly & $(168,24)$ & $(336,48)$ & $(504,168)$ \\
daily & $(7,1)$ & $(14,2)$ & $(30,7)$ \\
15-minute & $(96,4)$ & $(192,8)$ & $(672,96)$ \\
\bottomrule
\end{tabular}
\end{center}

Focused DGX fronts under \path{slurm/dgx/ablations/} change only the listed
values; their Selena counterparts are under \path{slurm/selena/ablations/}:
\begin{center}
\scriptsize
\begin{tabularx}{\linewidth}{@{}p{.30\linewidth}X@{}}
\toprule
Front & Values \\
\midrule
\path{ablation_n_datastore_dates.slurm} & test: $\{25,50,100\}$; publication:
$\{50,100,200,500\}$ \\
\path{ablation_n_fit.slurm} & fixed checkpoint value $\{10\}$ \\
\path{ablation_fit_stride.slurm} & $\Sfit\in\{1,P\}$ \\
\path{ablation_alpha.slurm} & fixed $\alpha\in\{0,10^{-4},10^{-3},10^{-2},10^{-1},1\}$; select $K$ \\
\path{ablation_k.slurm} & fixed $K\in\{1,3,5,10,15,20\}$; select $\alpha$ \\
\path{ablation_l.slurm} & $H=24$, $L\in\{24,96,168,336,504,512,720\}$ \\
\path{ablation_h.slurm} & $L=504$, $H\in\{24,48,64,96,168,336,504\}$ \\
\path{ablation_feature_design.slurm} & the seven designs in the signal table \\
\path{ablation_formulation.slurm} & ridge, delta ridge, and convex; shared and horizon forms \\
\path{ablation_fixed_protocol.slurm} & fixed-datastore crossed with fixed-training-set on one T3 grid \\
\path{ablation_general_scope.slurm} & all/same/other-user retrieval crossed with all/same-user fitting \\
\path{ablation_homogeneous.slurm} & all channels versus declared same-quantity channels on ETT and Weather \\
\path{ablation_backbones.slurm} & Chronos-2, Chronos-Bolt, Chronos-T5, and TS-ICL at $(512,64)$ \\
\path{ablation_sota_chronos_bolt.slurm} & causal Chronos-Bolt at $(512,64)$ beside published TS-RAG/Cross-RAG rows; protocols are labeled non-comparable \\
\bottomrule
\end{tabularx}
\end{center}

The Bayes baseline is part of the main adaptive workflow. CatBoost remains an
implemented supporting method but has no primary Slurm front.

\section{Code}

\subsection{Data objects}

Let $W=T-L-H+1$ be the number of valid window dates, $R$ the number of unique
query and fitting dates that require retrieval, $B$ the number of unique
window IDs that require a vanilla forecast, and $J$ the retrieval-feature
dimension. The source CSV is loaded once as \texttt{CsvTimeSeries}; overlapping
lookbacks and targets remain sliding views of $Z$.

\begin{center}
\scriptsize
\begin{tabularx}{\linewidth}{@{}p{.27\linewidth}p{.18\linewidth}X@{}}
\toprule
Object & Shape or key & Stored value and code owner \\
\midrule
\texttt{tables/windows} & $W$ dates, $U$ users & dates, timestamps, user names;
\path{src/proposal/extraction.py::_window_metadata} \\
\texttt{tables/window\_statistics} & $W\times U$ & lookback mean, standard
deviation, constant flag; same owner \\
\texttt{tables/window\_computation} & optional $A\times U\times J$ and
$B\times H$ & compact retrieval representations and forecasts, where $A$ is
the number of dates requiring a representation;
\path{src/pipeline/extraction.py::extract_online_features} \\
\texttt{tables/retrieval\_neighbors} & $R\times U\times\Kmax$ & neighbor
window ID and configured/raw/normalized distances; also $R$ dates, query flags,
and candidate counts; same owner \\
\path{context_cache/K_*/window_*.npy} & $(r,K)\mapsto\mathbb R^{U\times H}$
& requested context forecasts only;
\path{src/proposal/extraction.py::ContextForecastCache} \\
\path{selected_hyperparameters.csv} & $(q,u)$ or $q$ & selected $\alpha$,
$K$, validation loss, parameter count;
\path{src/pipeline/adaptation.py::evaluate_online_linear} \\
\path{coefficient_trajectory.npy} & query, optional user, optional horizon,
feature & refitted ridge coefficients; same owner \\
\path{per_user_date_metrics.csv} & $(q,u,\text{method})$ & (18)--(21)
and vanilla values; \path{src/proposal/ridge.py} \\
\path{metrics.json}, \path{compute_timing.json} & one result & aggregate
measures and complete/cold times; \path{src/results/efficiency.py} and
\path{src/pipeline/adaptation.py} \\
\bottomrule
\end{tabularx}
\end{center}

A stable window ID is
\[
 \operatorname{id}(r,u)=\operatorname{pos}(r)U+(u-1),         \tag{22}
\]
Here $\operatorname{pos}(r)$ is the zero-based position of date $r$ in the
valid-date array, so $0\le\operatorname{pos}(r)<W$.
\texttt{open\_extraction\_arrays} in
\path{src/pipeline/contracts.py} reconstructs source windows, retrieved values,
and sparse forecast joins from these IDs.

\subsection{Pipeline}

\begin{algorithm}[H]
\caption{Extraction}
\footnotesize
\begin{algorithmic}[1]
\Require panel $Z$; settings $L,H,M,\Nfit,\Sstore,\Sfit,\Sq,\Kmax$
\State build sliding views $x_{r,u},y_{r,u}$ and store their IDs, means, and scales
\State precompute T0, T1+T2, and every T3 query date
\State build $\mathcal Q$ from (6)
\ForAll{$q\in\mathcal Q$}
  \State build the fixed or rolling $\mathcal F_q$ from (5)
\EndFor
\State $\mathcal R\gets\mathcal Q\cup\bigcup_{q\in\mathcal Q}\mathcal F_q$
\ForAll{$r\in\mathcal F_{q_0}\cup\{q_0\}$ first, then remaining $r\in\mathcal R$ by date}
  \State build the fixed-T0 or rolling $\mathcal D_r$ from (3)--(4)
  \For{$u=1$ to $U$}
    \State exact-search the allowed users and save ranks $1{:}\Kmax$
  \EndFor
\EndFor
\State forecast each query or selected-neighbor window ID once
\State save extraction arrays, diagnostics, manifest, and time
\end{algorithmic}
\end{algorithm}

\begin{algorithm}[H]
\caption{Adaptation and evaluation}
\footnotesize
\begin{algorithmic}[1]
\Require extraction arrays; method; fitting scope; candidate $\alpha$ and $K$ values
\ForAll{$q\in\mathcal Q$ in chronological order}
  \State reconstruct $\mathcal F_q$ from (5)
  \If{method is ridge, delta ridge, or convex}
    \State build the chosen design from $V,C,\bar Y,Y_{1:K},N_{1:K}$
    \State select $\alpha,K$ on the newest validation dates when enabled
    \State refit on all $\mathcal F_q$ with (11) or (13)
    \State predict with (12) or $D\widehat\beta$
  \ElsIf{method is a gate}
    \State fit the mean-advantage or CatBoost score from (15)
    \State predict with (17)
  \ElsIf{method is direct covariate prediction}
    \State use $C_{q,u,K}$ without fitted parameters
  \Else
    \State run TS-RAG ARM on its raw retrieved sequences
  \EndIf
  \State save (18)--(19) for every user and the all-user date batch
\EndFor
\State aggregate (20)--(21), user measures, coefficients/importances, and time
\end{algorithmic}
\end{algorithm}

The direct code path is
\begin{center}
\small
\path{src/scripts/extract.py} $\longrightarrow$
\path{src/pipeline/extraction.py} $\longrightarrow$
\path{src/scripts/adapt.py} $\longrightarrow$
\path{src/pipeline/adaptation.py} $\longrightarrow$
\path{src/scripts/tables.py}.
\end{center}
Date-plan construction is in \path{src/proposal/date_planning.py} and datastore
selection is in \path{src/proposal/datastore.py}; ridge designs and solves
are in \path{src/proposal/ridge.py}; gates are in
\path{src/proposal/gates.py}; TS-RAG is in
\path{src/external_models/tsrag/} and \path{src/pipeline/tsrag.py}; run lists are
defined only in \path{src/pipeline/profiles.py}.

\section{Execution}

Each DGX front under \path{slurm/dgx/} has a matching
\texttt{\_selena.slurm} overflow front under \path{slurm/selena/}. Both set
the same family and call
\path{src/slurm/run_family.sh}. The default order is
\texttt{extract,adapt,tables}; \texttt{STAGES} may contain a recovery subset.
Each stage uses Hydra and the same seed. Completed runs with the exact recorded
settings are reused. Training and full inference run only on the cluster.
DGX arrays, metrics, figures, and tables are under \path{outputs/}, and run
logs are under \path{logs/}. Selena fronts preserve the same relative artifact
identity under \path{outputs_selena/} and write their streams under
\path{logs_selena/}, so returned overflow work never collides with DGX trees.
Selena uses partition \texttt{an} with QoS \texttt{an\_preemptable}.

Lightweight synchronization and publication transfer requested logs plus only
aggregate \path{outputs*/reports/} artifacts; they do not traverse per-run or
diagnostics trees. Detailed transfer adds non-binary outputs and diagnostics,
while full synchronization is reserved for binary recovery payloads.

No online-adaptation result has yet been analyzed. This document therefore
defines computations and runs only; it makes no performance claim.

\end{document}
